\section{Running all recolorings of the four-state graph}
\label{app:four-state-recolorings}
Start with $R=[4312]$, $B=[3443]$ and right-action composition.
Then
\[
 A=\frac12\begin{pmatrix}0&0&1&1\\0&0&1&1\\1&0&0&1\\0&1&1&0\end{pmatrix},
 \qquad
 \Delta=\frac12\begin{pmatrix}0&0&-1&1\\0&0&1&-1\\1&0&0&-1\\0&1&-1&0\end{pmatrix},
 \qquad \pi=\frac16(1,1,2,2).
\]
For a subset $S$ of source vertices, swap the red and blue outgoing
edges precisely at $S$. Writing $G_S=\operatorname{diag}(\epsilon_1,
\ldots,\epsilon_4)$, with $\epsilon_x=-1$ at swapped sources and $+1$
elsewhere, gives
\[
 A_S=A,\qquad\Delta_S=G_S\Delta,\qquad
 R_S=A+\Delta_S,\quad B_S=A-\Delta_S.
\]
The stationary row is unchanged. In particular $\operatorname{rank}A=3$
and $d=\operatorname{rank}\Delta_S=3$ for every recoloring.


\subsection*{Fixed target-pair structure}
Using Definition~\ref{def:target-pair-graph}, the compact source-indexed
list and the underlying undirected graph are
\[
 \boxed{\mathcal T(A)=(34,34,14,23),\qquad H_A:\ 1-4-3-2.}
\]
The two occurrences of $34$ come from sources 1 and 2. They represent a
repeated auxiliary edge, not parallel edges in the original directed graph.
This fixed header is recorded once, not repeated in each recoloring row.

The target-pair graph has one connected component, and it is bipartite,
with sides $\{1,3\}$ and $\{2,4\}$. Consequently
\[
 \ker\Delta_S=\operatorname{span}\{(1,1,1,1)^{\mathsf T}\},\qquad
 \ker A=\operatorname{span}\{(1,-1,1,-1)^{\mathsf T}\}.
\]
Thus $\operatorname{rank}\Delta_S=\operatorname{rank}A=3$ for all
sixteen recolorings. There are exactly two labeled CC choices, interchanged
by a global color swap: direct every red edge into $\{1,3\}$ and every
blue edge into $\{2,4\}$, or reverse these target types. They are rows
$k=14$ and $k=4$ of the table, respectively.

The graph-intrinsic canonical component partition is the single block
$\{1,2,3,4\}$, whereas its chosen bipartition has two sides. Neither is
to be identified automatically with a kernel partition. Most importantly,
this same auxiliary graph supports kernel ranks one, two and three.
It determines the ranks of $A$ and $\Delta$, and CC existence, but not
the kernel rank or whether $\pi$ generates the full row space of
$\Delta_S$ through positive powers. The latter is exactly what the
ranked test in the recoloring table measures.

\subsection*{A one-at-a-time enumeration scheme}
For $k=0,\ldots,15$, use the binary-reflected Gray code
$g_k=k\mathbin{\mathrm{xor}}\lfloor k/2\rfloor$.
Bit $x-1$ specifies whether source $x$ is swapped relative to the original
coloring. Consecutive masks differ in exactly one bit, so consecutive
colorings differ by a single elementary recoloring. Every labeled coloring
occurs once; global red/blue interchanges are retained, not identified.

For each coloring compute the cycle matrix
\[
 \Sigma_S=\begin{pmatrix}\pi\Delta_S\\\pi\Delta_S^2\\\pi\Delta_S^3\end{pmatrix}.
\]
The coloring is called ranked here precisely when
$\operatorname{rank}\Sigma_S=d$. The sequence starts with exponent one;
$\pi$ is not an additional row of this test.
Independently enumerate the transformations generated by $R_S,B_S$;
the smallest image size is the kernel rank $r_K$.
A breadth-first enumeration also gives shortest reset words.

\begin{table}[tbp]
\centering\small
\begin{tabular}{r l c c c c c}
\hline
$k$ & Swapped sources $S$ & $R_S$ & $B_S$ & $r_K$ & $\operatorname{rank}\Sigma_S$ & Reset word\\
\hline
0 & $\varnothing$ & $[4312]$ & $[3443]$ & 2 & 1 & $\text{--}$ \\
1 & $\{1\}$ & $[3312]$ & $[4443]$ & 1 & 3 & $RB$ \\
2 & $\{1,2\}$ & $[3412]$ & $[4343]$ & 2 & 1 & $\text{--}$ \\
3 & $\{2\}$ & $[4412]$ & $[3343]$ & 1 & 3 & $RB$ \\
4 & $\{2,3\}$ & $[4442]$ & $[3313]$ & 1 & 3 & $RB$ \\
5 & $\{1,2,3\}$ & $[3442]$ & $[4313]$ & 1 & 3 & $RRB$ \\
6 & $\{1,3\}$ & $[3342]$ & $[4413]$ & 3 & 1 & $\text{--}$ \\
7 & $\{3\}$ & $[4342]$ & $[3413]$ & 1 & 3 & $BBR$ \\
8 & $\{3,4\}$ & $[4343]$ & $[3412]$ & 2 & 1 & $\text{--}$ \\
9 & $\{1,3,4\}$ & $[3343]$ & $[4412]$ & 1 & 3 & $BR$ \\
10 & $\{1,2,3,4\}$ & $[3443]$ & $[4312]$ & 2 & 1 & $\text{--}$ \\
11 & $\{2,3,4\}$ & $[4443]$ & $[3312]$ & 1 & 3 & $BR$ \\
12 & $\{2,4\}$ & $[4413]$ & $[3342]$ & 3 & 1 & $\text{--}$ \\
13 & $\{1,2,4\}$ & $[3413]$ & $[4342]$ & 1 & 3 & $RRB$ \\
14 & $\{1,4\}$ & $[3313]$ & $[4442]$ & 1 & 3 & $RB$ \\
15 & $\{4\}$ & $[4313]$ & $[3442]$ & 1 & 3 & $BBR$ \\

\hline
\end{tabular}
\caption{All sixteen labeled recolorings in Gray-code order. Swapped
sources are measured relative to the initial coloring. Words act from
left to right. A dash means that no reset word exists. Every displayed
reset word has minimum length.}
\label{tab:four-state-all-recolorings}
\end{table}

There are ten rank-one kernels, four rank-two kernels, and two rank-three
kernels. In this example, exactly the ten rank-one colorings are ranked.
Thus the sufficient ranked criterion detects every synchronizing coloring
in this table; this computation does not establish a general converse.
The two CC colorings occur at $k=4$ and $k=14$; their global color
interchange exchanges them. Complementary swap sets always give a global
color interchange, hence the same kernel and cycle-matrix ranks.
The complete exact records, including kernel cardinalities, are supplied
in the accompanying JSON and CSV files.
